
Machine learning reproducibility failures waste researcher time when discovered during training. A researcher adapting published code may discover tensor shape mismatches, dtype inconsistencies, or cross-library device placement errors only after initiating costly training runs. These failures originate from environment-stage API configuration but surface during execution when mitigation is expensive.

Existing reproducibility practices address different failure modes. \textbf{Environment isolation} (Docker, virtual environments) prevents system-level conflicts. \textbf{Dependency pinning} (requirements.txt, conda) stabilizes library versions across machines. \textbf{Integration testing} (pytest) validates repository-specific logic. However, none validate library-level behavioral assumptions---the invariants that documented APIs should satisfy but may violate under version updates or environmental variations.

We investigate whether \textbf{executable API contracts} can shift defect detection earlier in the development lifecycle. API contracts encode three classes of invariants: (1) structural contracts validate tensor shapes, dtypes, and devices at import time, (2) metamorphic contracts enforce mathematical properties through lightweight runtime probes, and (3) composition contracts validate cross-library consistency.

This work presents proof-of-concept implementations and controlled validation experiments for each contract type. We make the following contributions:

\textbf{Structural Contract PoC}: We implement decorator-based structural contracts and validate on 2 real scenarios (shape mismatch, dtype mismatch), achieving 100\% detection at import time with <0.03s overhead.

\textbf{Metamorphic Contract PoC}: We implement softmax sum and dropout identity validators and evaluate on 50 synthetic scenarios, achieving 100\% detection (40/40 violations) with 0\% false positives (0/10 control cases) in 3.7ms average execution time.

\textbf{Cross-Framework Validation PoC}: We implement dtype preservation, shape consistency, and numerical tolerance contracts for cross-framework model conversion, achieving 100\% detection on 30 synthetic test cases (24 defects, 6 valid) with 0\% false positive rate and <1ms validation time.

These proof-of-concept results establish technical feasibility. Section 2 positions our work within related literature. Section 3 details our three-tier contract architecture. Section 4 describes experimental methodology. Section 5 presents validation results. Section 6 discusses limitations and scope boundaries. Section 7 concludes with future research directions.
